\section{Supporting subsystems and extensions}
\label{sec:ext}

These subsystems were evaluated separately from the joint-control study, on the same shared, GPU-less
host (B-001, B-005).

\subsection{Crash-consistent checkpoint engine}
\label{sec:ext-ckpt}

\texttt{aegis-ckpt} is a Rust storage server implementing P1 (\texttt{docs/protocols.md}): log-structured
segments whose records carry a BLAKE3 payload digest and a CRC32C header check, a RocksDB chunk index
written only after the segment fsync, a checksummed manifest WAL with group commit (a commit is acknowledged
only once durable), and mark/copy/sweep GC with epochs. Writes use io\_uring or \texttt{pwrite}. The Python
client snapshots tensors (the training stall), flushes in the background, chunks each tensor at a fixed
size (layouts never shift between steps, so no content-defined chunking; D-016), uploads only missing
chunks and verifies every chunk on restore.

\paragraph{Crash consistency.} A trial is \emph{ok} if the recovered latest checkpoint is at least the last
acknowledged one and every retained checkpoint restores bit-exactly, \emph{detected} if recovery or restore
raises an integrity, rollback or missing-chunk error, and \emph{undetected} otherwise.
(i)~\CckptKillTrials{} \texttt{kill -9} trials of the server at random times during save or GC
(\CckptKillSaves{} saves, \CckptKillGcs{} GC cycles, io\_uring and pwrite alternating): all
\CckptKillOk{} ok. The latest checkpoint was the last acknowledged one in \CckptKillLatestAcked{} trials
and a durable but unacknowledged newer one in \CckptKillLatestUnacked.
(ii)~\CckptTornTrials{} truncations of a random store file at a random offset: \CckptTornDetected{}
detected, \CckptTornOk{} ok, \CckptTornUndetected{} undetected. Every WAL and segment truncation was
detected. (a cut WAL loses acknowledged commits, which restore reports via a trusted sequence
high-water mark); index damage was repaired.
(iii)~\CckptFlipTrials{} trials flipping \cfg{1}--\cfg{3} random bits: \CckptFlipDetected{} detected,
\CckptFlipOk{} ok, \CckptFlipUndetected{} undetected; the \CckptFlipSegmentOk{} undetected segment flips
hit bytes that no retained checkpoint references.

\paragraph{Performance.} Table~\ref{tab:ckpt} compares the engine with \texttt{torch.save} and PyTorch
DCP on XFS for a small TinyVLM with AdamW state (\CckptStateFullMB\,MB; \CckptStateFrozenMB\,MB with frozen
embeddings and lower blocks). The async engine cuts the stall to the snapshot copy; a synchronous engine
save stalls \CckptSyncRatioMin--\CckptSyncRatioMax$\times$ longer. The engine is \emph{not} faster to
durability or recovery at this size: client-side hashing, encryption and a TCP hop cost more than a
sequential write of a few tens of MB, and training steps run \CckptOverlapMin--\CckptOverlapMax$\times$
slower during the background flush, which competes for the same \cfg{4} CPUs. On ext4 the stall CIs of engine and
\texttt{torch.save} overlap. Across
successive checkpoints of a larger preset (\CckptDedupStateMB\,MB), the logical/uploaded ratio is
\CckptDedupFull{} for full fine-tuning (AdamW rewrites every parameter and moment), \CckptDedupFrozen{}
with the lower half frozen and \CckptDedupAdapter{} when only the last block trains
(\CckptNewFracAdapter\% new bytes per checkpoint), independent of chunk size from \cfg{64}\,KiB to
\cfg{4}\,MiB.

\begin{table}[t]
\centering\footnotesize\setlength{\tabcolsep}{3pt}
\caption{Checkpointing on XFS (mean of \CckptBenchN{} trials; seconds unless stated). ur/pw: engine,
async, io\_uring/pwrite; sync: io\_uring, synchronous. Engine columns use P2 encryption.}
\label{tab:ckpt}
\begin{tabular}{@{}lrrrrr@{}}
\toprule
 & torch.save & DCP & ur & pw & sync \\
\midrule
stall, full & \CckptStallXfsFullTorch & \CckptStallXfsFullDcp & \CckptStallXfsFullUring & \CckptStallXfsFullPwrite & \CckptStallXfsFullSync \\
stall, frozen & \CckptStallXfsFrozenTorch & \CckptStallXfsFrozenDcp & \CckptStallXfsFrozenUring & \CckptStallXfsFrozenPwrite & \CckptStallXfsFrozenSync \\
durable, full & \CckptDurableXfsFullTorch & \CckptDurableXfsFullDcp & \CckptDurableXfsFullUring & \CckptDurableXfsFullPwrite & \CckptDurableXfsFullSync \\
recovery, full & \CckptRecoveryXfsFullTorch & \CckptRecoveryXfsFullDcp & \CckptRecoveryXfsFullUring & \CckptRecoveryXfsFullPwrite & \CckptRecoveryXfsFullSync \\
dev./logical, frozen & \CckptWampXfsFrozenTorch & \CckptWampXfsFrozenDcp & \CckptWampXfsFrozenUring & \CckptWampXfsFrozenPwrite & \CckptWampXfsFrozenSync \\
\bottomrule
\end{tabular}
\end{table}

\paragraph{Limits.} \texttt{kill -9} spares the page cache, no campaign models a disk that lies about
fsync, and torn writes hit a quiescent store. The filesystems are loop-mounted images on the Docker Desktop
VM disk (B-014). CIs are wide under shared load, ext4 and XFS ran at different loads, and at tens of MB
fixed per-checkpoint costs dominate.

\subsection{Encrypted checkpoints and certificate rotation}
\label{sec:ext-sec}

\paragraph{Envelope (P2).} Chunks are sealed with a per-job, per-epoch AES-256-GCM key wrapped by a
key-encryption key; chunk ids are keyed BLAKE3, so dedup works without exposing plaintext hashes; manifests
are Ed25519-signed and checked against a sequence high-water mark. On a \CenvStateMB\,MB checkpoint (\CenvTrials{} interleaved trials) a cold restore takes
\CenvColdPtwo\,s with P2 against \CenvColdPlain\,s in plaintext (\CenvColdRatio$\times$; Welch
$p=\CenvColdP$, Mann-Whitney $p=\CenvColdPmwu$). The difference, about \CenvColdAbs\,s, sits in the
verify/decrypt phase (\CenvVerifyPlain\,s$\to$\CenvVerifyPtwo\,s, Welch $p=\CenvVerifyP$). Encryption adds
\CenvSaveEncrypt\,s per save. The anti-rollback mark still lives in the client, not the controller as P2
specifies.

\paragraph{Rotation (P3).} Workload certificates (ECDSA P-256, SPIFFE URI SANs) have lifetime $T$ and a
\texttt{notBefore} backdated by $2S$, where $S$ bounds clock skew. Workloads renew at the midpoint of the
current certificate's validity and retry on failure. Servers reload credentials per handshake; clients open
a new channel on renewal. CA rotation has two phases: the new CA reaches every verifier's bundle before it
signs, and the old CA leaves only after all its certificates have expired (plus $S$). Both the backdating
and this ordering came from TLC counterexamples (\S\ref{sec:ext-formal}).
% TODO(main): rotation results
Results of the live rotation campaign under injected skew are pending.

\subsection{Formal checks}
\label{sec:ext-formal}

TLA+~\citep{tla} models of P1 and P3 are checked with TLC. Fixed protocols must pass; buggy variants and
non-vacuity checks must produce counterexamples. All \CtlcModels{} models behave as expected
(\CtlcPassing{} pass, \CtlcViolating{} counterexamples; longest run \CtlcMaxWall\,s).

\paragraph{P1.} Uploads with dedup, group commit, COMMIT, WAL fsync, crashes anywhere (any subset of
unsynced chunks and prefix of unsynced WAL survives), recovery and concurrent GC. Invariants: committed
manifests reference only durable chunks (I1), recovery returns the latest commit (I2), GC never removes a
chunk of a retained or accepted commit (I3), acknowledged commits survive, and liveness. Bounds: \cfg{3} chunks,
\cfg{2} manifests sharing one, $K=1$, one crash and one GC (\CtlcPoneFixedStates{} distinct states), and
two of each for safety (\CtlcPoneBigStates). As first written, GC snapshotted the live set at mark time.
TLC found that a commit accepted but not yet fsync'd was swept (I3, \CtlcPoneAsWrittenTrace-state trace), and
that a commit landing between mark and sweep lost its deduplicated chunk, so that after the WAL fsync a
committed manifest referenced a deleted chunk (I1, \CtlcPoneAsWrittenCommittedTrace{} states). P1 and the
engine now refresh the epoch of every referenced chunk at COMMIT under the WAL lock, treat unsynced WAL
records as live, and re-check epochs at sweep time.

\paragraph{P3.} Discrete time, per-node clock offsets in $[-S,S]$, renewal with failures and back-off, and a
full CA rotation. R1: every verifier
accepts the certificate the workload presents, in every state. Bounds: $T=\cfg{10}$, $S=1$, one failed
issuance, a \cfg{26}-tick horizon (\CtlcPthreeFixedStates{} states). P3 originally set \texttt{notBefore}
to $\textit{now}-S$; TLC found a \CtlcPthreeAsWrittenTrace-state trace in which a CA running $S$ fast
issues a certificate that a verifier running $S$ slow rejects. Some held certificate is always still
accepted, but workloads present their newest one. Backdating by $2S$ fixes it and tightens the
precondition from retries $<T/2-S$ to $<T/2-2S$; with two failed issuances TLC finds a violation at
$T=\cfg{10}$ and none at $T=\cfg{16}$. Seeded bugs (late renewal, early CA removal, signing before distribution;
index-before-fsync and epoch-blind sweep in P1) each yield a counterexample.

\paragraph{Limits.} Bounds are small; the models assume atomic chunk writes, perfect checksums, no clock
drift and one workload; and TLC checks the models, not the code.

\subsection{Real-time media}
\label{sec:ext-media}

A live assistant session streams a \cfg{256}$\times$\cfg{256} chart screen-share as AV1 over our own
RTP/RTCP stack with Opus audio. Speech goes to
whisper-tiny; vision is a stub (B-010). The
\emph{independent} policy uses a GCC-style rate controller, fixed frame rate, fixed-rate frame sampling and
a FIFO inference queue. \emph{Joint} also picks the encoder frame rate from an offline quality table, caps
sampling at \cfg{0.8} of the measured inference rate, and uses a latest-wins vision slot that speech jobs
can jump.

\paragraph{Virtual time.} The campaign runs on a virtual clock (D-013): codec, STT and vision really
execute and are charged their measured CPU time, and the network is a seeded in-process link model, not
kernel netem. All latencies below are therefore \emph{model-derived}, not wall-clock; wall-clock timing
was meaningless on this host (load average \CmediaLoadMin--\CmediaLoadMax{} on \cfg{18} cores).

\paragraph{Results (mixed).} \CmediaTrials{} trials: \CmediaConditions{} conditions $\times$
\CmediaSeeds{} seeds. When inference is the bottleneck (vision service time \cfg{150} instead of
\cfg{50}\,ms), joint answers in \CmediaOverloadGaJoint\,s against \CmediaOverloadGaIndep\,s (Welch
$p=\CmediaOverP$), and backpressure drops fall from \CmediaOverDropIndep{} to \CmediaOverDropJoint{} per
session. After Holm correction over \CmediaTests{} tests only the drop count stays significant
($p_\text{Holm}=\CmediaOverDropHolm$; latency $p_\text{Holm}=\CmediaOverHolm$), and part of the gain is the
baseline's fixed sampling rate, which the stub vision cannot sustain. Under a bandwidth collapse joint
answers in \CmediaBwdropGaJoint\,s against \CmediaBwdropGaIndep\,s ($p=\CmediaBwdropP$,
$p_\text{Holm}=\CmediaBwdropHolm$); VMAF (\CmediaBwdropVmafJoint{} vs \CmediaBwdropVmafIndep) and stalls
point the same way but are not significant. Under loss and jitter there is no significant difference, and
in the clean profile joint is slightly slower (\CmediaCleanGaJoint{} vs \CmediaCleanGaIndep\,s). The
encoder knob was inert: the quality table rates \cfg{30}\,fps highest at every bitrate (e.g.\ VMAF
\CmediaCalFourThirty{} vs \CmediaCalFourTen{} at \cfg{10}\,fps, \cfg{400}\,kbps), so joint always chose it.
Blocking the media path on inference (an ablation) is far worse: \CmediaOverloadStallBlock\,s of stalls
per session and VMAF \CmediaOverloadVmafBlock. (A reordering jitter profile triggered a NACK storm under both
policies; our receiver has no reordering tolerance, so we report order-preserving jitter.)

\paragraph{WebRTC leg and limits.} aiortc lacks AV1 RTP (B-008), so a separate aiortc VP8 ingest leg with
server-side AV1 transcode was measured on localhost, unimpaired, in wall time: \CrtcGsMean\,ms mean
(\CrtcGsP\,ms p95) glass-to-server, \CrtcDelivery\% of frames delivered, \CrtcTranscode\,ms per frame to
transcode. The network is emulated, vision is a stub, latencies are model outputs, and power is low.

\subsection{Kernels}
\label{sec:ext-kern}

No NVIDIA GPU was available (B-001). We wrote a fused, flash-style Triton~\citep{triton,flashattn} kernel
for variable-length prefix-LM attention over packed sequences (bidirectional image prefix, causal text, no
cross-sequence attention) and a Metal port of the same algorithm. Against an fp64 reference with max-error
tolerances $\CkernTolSingle$ (fp32) and $\CkernTolHalf$ (fp16), Triton passes
\CkernTritonPass/\CkernTritonCases{} cases (worst errors $\CkernTritonErrSingle$, $\CkernTritonErrHalf$)
and Metal \CkernMetalPass/\CkernMetalCases{} ($\CkernMetalErrSingle$, $\CkernMetalErrHalf$). The Triton runs
use \texttt{TRITON\_INTERPRET=1}, a NumPy emulator: they check indexing and numerics, not GPU execution.

On the Apple GPU at the model's shape (fp32), the Metal kernel takes \CkernMpsModelSingleMetal\,ms against
\CkernMpsModelSingleSdpa\,ms for SDPA (overlapping CIs) and \CkernMpsModelSingleEager\,ms for padded eager
attention. At a larger shape it is \CkernMpsLargeSingleMetalOverSdpa$\times$ (fp32) and
\CkernMpsLargeHalfMetalOverSdpa$\times$ (fp16) \emph{slower} than SDPA, beating only eager; it uses no
simdgroup matrix instructions, a likely but untested cause. For CUDA we compiled the Triton kernel offline for
sm80 and inspected the PTX: with divisibility hints, fp16 code uses \CkernIrHalfHintsCpAsync{}
\texttt{cp.async} copies (none without), \CkernIrHalfHintsLdmatrix{} \texttt{ldmatrix} (vs
\CkernIrHalfPlainLdmatrix) and \CkernIrHalfHintsMma{} \texttt{mma} instructions. This shows the compiler
emits asynchronous tensor-core code; it says nothing about speed on real hardware.

\subsection{Post-training and privacy}
\label{sec:ext-pt}

All runs use synthetic data and our own char-level transformer (\CptParams{} parameters), as no usable
pretrained model or GPU was available (B-016). The task is tool-using question answering over a hidden
synthetic statement, scored by exact match. All methods start from the same per-seed warm start.

\paragraph{Equal compute.} Each method gets the same analytic FLOP budget after the warm start
(\cfg{2}$N$ per forward and \cfg{6}$N$ per trained token, sampling and scoring included). SFT trains on more
gold demonstrations; DPO~\citep{dpo} is online with best-vs-worst pairs under a verifiable reward;
GRPO~\citep{deepseekmath} uses the same reward; GRPO-RM uses a learned reward model. At $\CptLoBudget$
FLOPs exact match is \CptLoSftEm{} (SFT), \CptLoDpoEm{} (DPO), \CptLoGrpoEm{} (GRPO) and \CptLoGrpoRmEm{}
(GRPO-RM) against \CptLoWarmEm{} for the warm start; SFT is bimodal across seeds and not significant at
$n=\CptSeeds$ (Welch $p\le\CptLoSftPmax$). At $\CptHiBudget$ FLOPs SFT reaches \CptHiSftEm{} while DPO
(\CptHiDpoEm) and GRPO (\CptHiGrpoEm) stay at the warm-start level ($p<\CptHiSftPbound$). The RL methods
raised tool-call validity (DPO \CptLoDpoTool{} vs \CptLoWarmTool) but not correctness; the warm start never
solves three of six question types, so outcome rewards rarely see a success there (our interpretation). \emph{Reward hacking:} the reward model is accurate on held-out warm-start rollouts (AUC
\CptRmAuc), yet under GRPO-RM the training reward rises from \CptRmRewFirstA{} to \CptRmRewLastA{} (seed~0)
while exact match ends at \CptHiGrpoRmEm{} and tool-call validity collapses on three of five seeds
(\CptRmToolA, \CptRmToolC, \CptRmToolE; plain GRPO \CptHiGrpoTool).

\paragraph{DP-SGD and membership inference.} We fine-tune a public task model (exact match \CdpBaseEm) on
a small private set plus random-digit canaries with Opacus DP-SGD~\citep{dpsgd,opacus}: example-level,
$\delta=\cfg{10}^{\cfg{-5}}$, Poisson rate \CdpRate, \CdpSteps{} steps, RDP accountant with a PRV
cross-check. Attacks: a loss threshold~\citep{shokri} and offline LiRA with shadows trained on
same-format reference canaries (our first shadows saw no canaries and saturated LiRA at chance). Utility
is flat (Table~\ref{tab:dp}): every $\varepsilon$ from \CdpOneRdp{} to \CdpTenTwentyFourRdp{} stays at the
public model's level, while non-private fine-tuning reaches \CdpNoneEm{} and clipping without noise
\CdpClipEm, so the noise, not the clipping, removes utility in this small-data, short-schedule regime.
Leakage still rises with $\varepsilon$: large $\varepsilon$ buys leakage and no utility.

\begin{table}[t]
\centering\footnotesize\setlength{\tabcolsep}{3pt}
\caption{DP-SGD: privacy spent, exact match and membership-inference AUC (mean of \CdpSeeds{} seeds;
LiRA offline, pooled variance).}
\label{tab:dp}
\begin{tabular}{@{}lrrrrr@{}}
\toprule
$\varepsilon$ (RDP) & $\varepsilon$ (PRV) & $\sigma$ & exact & loss AUC & LiRA AUC \\
\midrule
$\infty$ (no DP) & -- & -- & \CdpNoneEm & \CdpNoneLossAuc & \CdpNoneLiraAuc \\
$\infty$ (clip only) & -- & 0 & \CdpClipEm & \CdpClipLossAuc & \CdpClipLiraAuc \\
\CdpTenTwentyFourRdp & \CdpTenTwentyFourPrv & \CdpTenTwentyFourSigma & \CdpTenTwentyFourEm & \CdpTenTwentyFourLossAuc & \CdpTenTwentyFourLiraAuc \\
\CdpTwoFiftySixRdp & \CdpTwoFiftySixPrv & \CdpTwoFiftySixSigma & \CdpTwoFiftySixEm & \CdpTwoFiftySixLossAuc & \CdpTwoFiftySixLiraAuc \\
\CdpSixtyFourRdp & \CdpSixtyFourPrv & \CdpSixtyFourSigma & \CdpSixtyFourEm & \CdpSixtyFourLossAuc & \CdpSixtyFourLiraAuc \\
\CdpSixteenRdp & \CdpSixteenPrv & \CdpSixteenSigma & \CdpSixteenEm & \CdpSixteenLossAuc & \CdpSixteenLiraAuc \\
\CdpFourRdp & \CdpFourPrv & \CdpFourSigma & \CdpFourEm & \CdpFourLossAuc & \CdpFourLiraAuc \\
\CdpOneRdp & \CdpOnePrv & \CdpOneSigma & \CdpOneEm & \CdpOneLossAuc & \CdpOneLiraAuc \\
\bottomrule
\end{tabular}
\end{table}

\paragraph{Federated.} Three in-process non-IID sites run FedAvg~\citep{fedavg} and FedProx for
\CfedRounds{} rounds of \CfedLocalSteps{} local steps, compared with centralised and local-only training at
equal SGD steps. FedAvg reaches \CfedAvgEm{} global exact match against \CfedCentralEm{} centralised
(\CfedAvgOverCentral\%, $p=\CfedCentralP$), below the best local-only site (\CfedLocalUsEm); FedProx does
not help (\CfedProxEm). It moves \CfedTotalMB\,MB of weights; the raw data are \CfedRawMB\,MB.

\paragraph{Limits.} The model is tiny and the task templated; nothing here transfers to foundation-model
scale, and the method ranking may change with scale, warm start, reward and budget. Compute matching is
analytic, and reward-model training is not charged to GRPO-RM. The DP guarantee covers only the private
fine-tuning step (Opacus secure mode off); few shadow models make the AUCs lower bounds on leakage.
Envelope encryption protects checkpoints, not what a model reveals through its outputs. Federated averaging
alone gives no formal privacy, and the simulation has no secure aggregation, stragglers or WAN.
